\documentclass[a4paper, 12pt, oneside]{book}

\usepackage[T1]{fontenc}
\usepackage[utf8]{inputenc}
\usepackage[spanish, es-lcroman]{babel}
\usepackage{float}
\usepackage{hyperref}
\usepackage{tocloft}
\usepackage{titling}
\usepackage{eurosym}
\usepackage{graphicx}
\usepackage{titlepic}
\usepackage{remreset}
\usepackage{bookmark}
\usepackage{multicol}

\renewcommand{\thechapter}{\Roman{chapter}}

\renewcommand{\labelitemi}{$\bullet$}
\renewcommand{\labelitemii}{$\circ$}
\renewcommand{\labelitemiii}{--}

\renewcommand{\theenumi}{\alph{enumi}}
\renewcommand{\labelenumi}{\theenumi)}

\renewcommand{\labelenumii}{\theenumii)}
\renewcommand{\theenumii}{\roman{enumii}}

\makeatletter
  \renewcommand{\thesection}{\@arabic\c@section}
  \@removefromreset{section}{chapter}
\makeatother

\setlength{\parskip}{1em}
\setcounter{tocdepth}{3}
\setcounter{secnumdepth}{3}

\hypersetup{
    colorlinks=true,
    linkcolor=black,
    filecolor=magenta,      
    urlcolor=cyan,
}

\hyphenation{Py-thon}

\title{
    \huge \textbf{Python Canarias} \\
    \textit{Estatutos de la asociación} \\
    \vspace{1cm}
    \large \textsc{Longa Vita Draconibus}
}
\titlepic{\includegraphics[width=.5\textwidth]{logo_python_canarias.png}}
\date{\textbf{1 de octubre de 2018}}
\author{}

\begin{document}

\maketitle

\tableofcontents

\chapter{De la asociación en general}

\section{La denominación}

De conformidad con el artículo 22 de la Constitución Española se constituye la entidad, no lucrativa, denominada ``Asociación cultural Python Canarias'', dotada de personalidad jurídica propia y capacidad de obrar y que se regirá por la Ley Orgánica 1/2002, de 22 de marzo, reguladora del Derecho de asociación, la Ley 4/2003, de 28 de febrero, de Asociaciones de Canarias, el Decreto 12/2007, de 5 de febrero, por el que se aprueba el Reglamento de Asociaciones de Canarias, demás disposiciones complementarias y por los presentes Estatutos.

\section{El domicilio social}

La asociación tendrá su domicilio social en la siguiente dirección postal:

\begin{quotation}

    Avda. Alonso Fernández de Lugo, 4 - 1A

    Edificio El Mayorazgo

    38300

    La Orotava

    Santa Cruz de Tenerife

    España

    % C/ La Higuera, 4 (Edificio Anaga) Portal 2B, primero

    % 38204

    % San Cristóbal de La Laguna

    % Santa Cruz de Tenerife

    % España
\end{quotation}

\section{El plazo de duración}

La duración de la constitución de la asociación será por tiempo indefinido.

\section{Finalidad}

Los fines de la asociación serán los siguientes:

\begin{enumerate}
    \item Promover el uso del lenguaje de programación Python, y de otros lenguajes por extensión.
    \item Servir como punto de encuentro de todas aquellas personas interesadas en el uso del lenguaje de programación Python, y de otros lenguajes por extensión.
    \item Dar soporte a quienes deseen aprender y desarrollar el uso del lenguaje de programación Python, y de otros lenguajes por extensión.
    \item Promover la programación como una competencia educativa básica.
\end{enumerate}

\section{Actividades}

Se desarrollarán toda clase de actividades que permitan a la asociación cumplir sus fines, principalmente dirigidos a la difusión del lenguaje de programación Python, tales como:

\begin{enumerate}
    \item Organizar cursos/talleres formativos.
    \item Participar en eventos.
    \item Vender artículos.
    \item Colaborar con otras organizaciones.
    \item Promover ofertas de empleo.
\end{enumerate}

\section{Ámbito territorial}

El ámbito territorial de la asociación en el que desarrollará sus actividades será la Comunidad Autónoma de Canarias.

\chapter{De los socios}

\section{Requisitos para ser socios}

Podrán ser miembros de la asociación:

\begin{enumerate}
    \item Las personas físicas, con capacidad de obrar, no sujetas a ninguna condición legal para el ejercicio del derecho de asociación, así como los/as menores de edad, con catorce o más años, siempre que cuenten con el consentimiento escrito y expreso de sus representantes legales.
    \item Las personas jurídicas.
\end{enumerate}

\section{Procedimiento para adquirir la condición de socio}

La condición de socio se adquirirá a solicitud de la persona interesada, manifestando su voluntad de contribuir al logro de los fines asociativos, dirigida al órgano de representación, y tras el visto bueno de éste.

\section{Pérdida de la condición de socio}

Se perderá la condición de socio por alguna de las situaciones siguientes:

\begin{enumerate}
    \item Por voluntad de la persona, manifestada al órgano de representación.
    \item Por acuerdo adoptado por el órgano de representación, conforme al régimen disciplinario establecido en estos estatutos.
    \item Por impago de la cuota de socio.
\end{enumerate}

\section{Clases de socios}

Los socios pueden ser:

\begin{enumerate}
    \item Fundadores/as: Quienes suscribieron el acta fundacional de la asociación.
    \item De número: Quienes hayan ingresado con posterioridad a la firma del acta fundacional de la asociación y se les admita como tales conforme a estos estatutos.
    \item Honoríficos/as: Quienes por decisión de la asamblea general colaboren de forma notable en el desarrollo de los fines de la asociación y/o quienes destaquen por ayudar con medios económicos y materiales a la asociación.
\end{enumerate}

\section{Derechos de los socios}

Los socios fundadores y los socios de número tendrán los derechos establecidos en el artículo 13 de la LEY 4/2003, de 28 de febrero, de Asociaciones de Canarias.

Los socios honoríficos tendrán los derechos establecidos en el artículo 13 de la LEY 4/2003, de 28 de febrero, de Asociaciones de Canarias, excepto el derecho a voto mencionado en el artículo 13 apartado a) de la citada ley.

\section{Obligaciones de los socios}

Los socios fundadores y los socios de número tendrán las obligaciones establecidas en el artículo 15 de la LEY 4/2003, de 28 de febrero, de Asociaciones de Canarias.

Los socios honoríficos tendrán las obligaciones establecidas en el artículo 15 de la LEY 4/2003, de 28 de febrero, de Asociaciones de Canarias, excepto el apartado b) del artículo 15 de la citada ley.

Los socios estarán obligados a notificar al órgano de representación cualquier modificación de sus datos personales aportados inicialmente.

\section{Publicidad de los datos}

Se podrá consultar la condición de socio a través de un formulario en la web de Python Canarias, donde se ingresará el DNI para su búsqueda.

\section{Notificaciones}

El medio de comunicación de notificaciones a los socios será vía correo electrónico y/o página web de la asociación.

\section{Participación de voluntarios}

El voluntariado es bienvenido para el desarrollo de las actividades de la asociación. Su participación en cada evento será establecida por el órgano de representación, o en su caso, por la asamblea general.

\chapter{Del régimen disciplinario}

\section{Infracciones}

\subsection{Infracciones muy graves}

Son infracciones punibles dentro del orden social y serán consideradas como muy graves:

\begin{enumerate}
    \item Todas aquellas actuaciones que perjudiquen u obstaculicen la consecución de los fines de la asociación, cuando tengan consideración de muy graves.  
    \item El incumplimiento o las conductas contrarias a las disposiciones estatutarias y/o reglamentarias de la asociación, cuando se consideren como muy graves.  
    \item El incumplimiento de los acuerdos válidamente adoptados por los órganos de la asociación, cuando se consideren muy graves.  
    \item La reiteración de una falta grave.  
    \item La protesta o actuaciones airadas y ofensivas que impidan la celebración de asambleas o reuniones del órgano de representación.  
    \item Participar, formular o escribir, mediante cualquier medio de comunicación social, manifestaciones que perjudiquen de forma muy grave la imagen de la asociación.  
    \item La usurpación ilegítima de atribuciones o competencias sin contar con la preceptiva autorización del órgano competente de la entidad.  
    \item Agredir, amenazar o insultar gravemente a cualquier socio.  
    \item La inducción o complicidad, plenamente probada, a cualquier socio en la comisión de las faltas contempladas como muy graves.  
    \item Cualquier acción o comentario que atente contra la igualdad efectiva entre hombres y mujeres.  
    \item El quebrantamiento de sanciones impuestas por falta grave o muy grave.  
    \item Todas las infracciones tipificadas como leves o graves y cuyas consecuencias físicas, morales o económicas, plenamente probadas, sean consideradas como muy graves.
\end{enumerate}

\subsection{Infracciones graves}

Son infracciones punibles dentro del orden social y serán consideradas como graves:

\begin{enumerate}
    \item El quebrantamiento de sanciones impuestas por infracciones leves.  
    \item Participar, formular o escribir mediante cualquier medio de comunicación social, manifestaciones que perjudiquen de forma grave la imagen de la asociación.  
    \item La inducción o complicidad, plenamente probada, de cualquier socio en la comisión de cualquiera de las faltas contempladas como graves.  
    \item Todas las infracciones tipificadas como leves y cuyas consecuencias físicas, morales o económicas, plenamente probadas, sean consideradas graves.  
    \item La reiteración de una falta leve.  
    \item El incumplimiento o las conductas contrarias a las disposiciones estatutarias y/o reglamentarias de la asociación, cuando se consideren como graves.  
    \item El incumplimiento de los acuerdos válidamente adoptados por los órganos de la asociación, cuando tengan la consideración de grave.
\end{enumerate}

\subsection{Infracciones leves}

Son infracciones punibles dentro del orden social y serán consideradas como leves:

\begin{enumerate}
    \item Todas aquellas conductas que impidan el correcto desarrollo de las actividades propias de la asociación, cuando tengan la consideración de leve.  
    \item El maltrato de los bienes muebles o inmuebles de la asociación.  
    \item Toda conducta incorrecta en las relaciones con los socios.  
    \item La inducción o complicidad, plenamente probada, de cualquier socio en la comisión de las faltas contempladas como leves.  
    \item El incumplimiento o las conductas contrarias a las disposiciones estatutarias y/o reglamentarias de la entidad, cuando se consideren como leves.
\end{enumerate}

\subsection{Infracciones de los miembros del órgano de representación}

\begin{enumerate}
    \item Se consideran infracciones muy graves susceptibles de ser cometidas por los miembros del órgano de representación:
    \begin{enumerate}
        \item La no convocatoria en los plazos y condiciones legales, de forma sistemática y reiterada, de los órganos de la asociación.  
        \item La incorrecta utilización de los fondos de la entidad.  
        \item El abuso de autoridad y la usurpación ilegítima de atribuciones o competencias.  
        \item La inactividad o dejación de funciones que suponga incumplimiento muy grave de sus deberes estatutarios y/o reglamentarios.  
        \item La falta de asistencia, en tres ocasiones y sin causa justificada, a las reuniones del órgano de representación. Se admitirá la asistencia y el voto por vía telemática siempre que se garanticen los medios necesarios por parte del órgano de representación.
    \end{enumerate}
    \item Se consideran infracciones graves:
    \begin{enumerate}
        \item No facilitar a los socios la documentación de la entidad que por éstos le sea requerida (estatutos, actas, normas de régimen interno, etc.).  
        \item No facilitar el acceso de los socios a la documentación de la entidad.  
        \item La inactividad o dejación de funciones cuando causen perjuicios de carácter grave al correcto funcionamiento de la entidad.
    \end{enumerate}
    \item Tienen la consideración de infracciones leves:
    \begin{enumerate}
        \item La inactividad o dejación de funciones, cuando no tengan la consideración de muy grave o grave.  
        \item La no convocatoria de los órganos de la asociación en los plazos y condiciones legales. 
        \item Las conductas o actuaciones contrarias al correcto funcionamiento del órgano de representación.  
        \item La falta de asistencia a una reunión del órgano de representación, sin causa justificada.
    \end{enumerate}
\end{enumerate}

\section{Sanciones}

Las sanciones susceptibles de aplicación por la comisión de infracciones muy graves serán la pérdida de la condición de socio o la suspensión temporal en tal condición durante un período de un año a cuatro años, en adecuada proporción a la infracción cometida.

Las infracciones graves darán lugar a la suspensión temporal en la condición de socio durante un período de un mes a un año, en adecuada proporción a la infracción cometida.

La comisión de las infracciones de carácter leve dará lugar a la amonestación, suspensión temporal del socio por un período de un mes y/o pérdida del derecho a voto del socio, esto último conforme al artículo 14 de la LEY 4/2003, de 28 de febrero, de Asociaciones de Canarias.

Las infracciones de los miembros del órgano de representación, en el caso de las muy graves al cese en sus funciones de miembro del órgano de representación y, en su caso, a la inhabilitación para ocupar nuevamente cargos en el órgano de gobierno; en el caso de las graves, el cese durante un período de un mes a un año, en adecuada proporción a la infracción cometida; y si la infracción cometida tiene el carácter de leve en la amonestación, suspensión temporal por el período de un mes y/o pérdida del derecho a voto del socio, esto último conforme al artículo 14 de la LEY 4/2003, de 28 de febrero, de Asociaciones de Canarias.

\subsection{Procedimiento sancionador}

Para la adopción de medidas disciplinarias acordes a la gravedad de la infracción cometida, se llevará una propuesta de sanción a la siguiente reunión del órgano de representación, que ratificará la misma mediante acuerdo.

A continuación, se notificará la sanción provisional al infractor, el cual dispone de 48 horas para alegaciones. En caso de haber alegaciones, se tratarán en la siguiente reunión del órgano de representación, al final de la cual se acordará una sanción definitiva que será notificada al infractor. En caso de no haber alegaciones, la sanción provisional pasaría a convertirse en definitiva.

\subsection{Prescripción}

Las infracciones prescribirán a los tres años, al año o al mes, según se trate de las muy graves, graves o leves, comenzándose a contar el plazo de prescripción al día siguiente a la comisión de la infracción.

\chapter{De los órganos sociales}

\section{Órgano de representación}

\subsection{Estructura del órgano de representación}

El órgano de representación será la junta directiva y estará integrada por los siguientes miembros:

\begin{enumerate}
    \item Un/a presidente/a.  
    \item Un/a vicepresidente/a.  
    \item Un/a secretario/a.  
    \item Un/a tesorero/a.  
    \item 4 vocales.      
\end{enumerate}

\subsection{Competencias del órgano de representación}

Las facultades del órgano de representación se extenderán, con carácter general, a todos los actos propios de las finalidades de la asociación, siempre que no requieran, conforme a los presentes estatutos, autorización expresa de la asamblea general.

En particular son facultades del órgano de representación:

\begin{enumerate}
    \item Velar por el cumplimiento de los estatutos y ejecutar los acuerdos tomados en las asambleas generales.  
    \item Confeccionar, custodiar y poner a disposición de los socios, si así se requiere: memorias, actas, contabilidad, inventario, registro de socios y presupuestos de la asociación.  
    \item Acordar la celebración de actividades.  
    \item Recaudar la cuota de los socios y administrar los fondos sociales.  
    \item Instruir procedimiento sancionador frente a infracciones de los socios que estén recogidas en los estatutos.  
    \item Proponer a la asamblea general para su aprobación las cuentas anuales y el presupuesto formulados por el tesorero, así como la memoria de actividades formuladas por el secretario.
\end{enumerate}

\subsection{Cargos del órgano de representación}

Sólo podrán formar parte del órgano de representación los socios.

Los cargos que constituyen el órgano de representación, que serán voluntarios, y por tanto deberán ser aceptados expresamente en documento firmado por las personas designadas, tendrán una duración de 4 años, pudiendo ser reelegidos.

Los miembros del órgano de representación comenzarán a ejercer sus funciones una vez aceptado el mandato para el que hayan sido designados por la asamblea general.

Las vacantes que se produzcan en el órgano de representación, antes de terminar su período de mandato, serán cubiertas por los socios que designe la propia junta directiva, dando cuenta de las sustituciones en la primera asamblea general que se celebre, debiendo ratificarse dicho acuerdo por la asamblea. En caso contrario, se procederá a la elección del socio que debe cubrir la vacante en la misma sesión de la asamblea.

\subsubsection{Facultades del presidente/a}

Serán facultades del presidente/a: 

\begin{enumerate}
    \item Ostentar la representación legal de la asociación. 
    \item Convocar y presidir las reuniones del órgano de representación y de la asamblea general de acuerdo con lo establecido en los presentes Estatutos. 
    \item Velar por el cumplimiento de los fines sociales. 
    \item Autorizar con su firma las actas, certificaciones y demás documentos de la Asociación. 
    \item Solicitar, percibir o hacer efectivas las ayudas, subvenciones o patrocinios que por cualquier concepto provengan de organismos públicos o entidades privadas, así como contratar y convenir con entes públicos de cualquier naturaleza y privados.
    \item Y cuantas facultades se le confieran, no expresamente asignadas a otros órganos.
\end{enumerate}

\subsubsection{Facultades del vicepresidente/a}

Serán facultades del vicepresidente/a:

\begin{enumerate}
    \item Sustituir al presidente/a en caso de vacante, ausencia o enfermedad asumiendo sus funciones, con carácter provisional, cuando el titular cesare en el cargo.
    \item Las que les deleguen el presidente/a o la asamblea general.
\end{enumerate}

\subsubsection{Facultades del secretario/a}

Serán facultades del secretario/a:

\begin{enumerate}
    \item Redactar y certificar las actas de las sesiones de las asambleas generales y del órgano de representación. 
    \item Llevar el libro del registro de asociados, consignando en ellos la fecha de su ingreso y las bajas que hubieren. 
    \item Recibir y tramitar las solicitudes de ingreso.
    \item Llevar una relación del inventario de la asociación.
    \item Tener bajo su custodia los documentos y archivos de la asociación. 
    \item Expedir certificaciones. 
    \item Formular la memoria de actividades.
\end{enumerate}

\subsubsection{Facultades del tesorero/a}

Serán facultades del tesorero/a:

\begin{enumerate}
    \item Tendrá a su cargo los fondos pertenecientes a la asociación.
    \item Elaborar los presupuestos, balances e inventarios de la asociación. 
    \item Firmará los recibos, cobrará las cuotas de los asociados y efectuará todos los cobros y pagos.
    \item Llevar y custodiar los libros de contabilidad. 
    \item Formular las cuentas anuales y el presupuesto.
\end{enumerate}

\subsubsection{Facultades de los/las vocales}

Los/as vocales tendrán el derecho y la obligación de asistir a las sesiones del órgano de representación con voz y con voto así como, en régimen de delegación, podrán desempeñar las funciones que les confiera el órgano de representación.

\subsection{Convocatoria del órgano de representación}

Las reuniones del órgano de representación se celebrarán previa convocatoria del presidente, acompañada del orden del día consignando lugar, fecha y hora. En el caso de que la convocatoria no incluyese el lugar de celebración se entenderá a todos los efectos el domicilio social.

La toma de decisiones se realizará por mayoría de votos de los socios asistentes.

\section{Asamblea General}

\subsection{Estructura de la asamblea general}

La asamblea general es el órgano supremo de la asociación, integrada por todos los asociados, que adoptará sus acuerdos bajo el principio democrático de mayorías de los votos válidamente emitidos.

\subsection{Competencias de la asamblea general}

Corresponde a la asamblea general, deliberar y tomar acuerdos sobre los siguientes asuntos:

\begin{enumerate}
    \item Examinar y aprobar el plan general de actuación y la memoria anual de actividades que le presente el órgano de representación. 
    \item Aprobar el presupuesto anual de gastos e ingresos del siguiente año y el estado de cuentas del ejercicio anterior. 
    \item Decidir sobre la disposición o enajenación de bienes. 
    \item Elegir y separar a los miembros de el órgano de representación. 
    \item Solicitar la declaración de utilidad pública o interés público canario.
    \item Acordar la integración de la asociación en Federaciones o Confederaciones, así como la separación de las mismas. 
    \item Controlar la actividad de el órgano de representación y aprobar su gestión. 
    \item Modificar los estatutos. 
    \item Acordar la disolución de la asociación. 
    \item Designar la comisión liquidadora en caso de disolución. 
    \item Acordar la remuneración de los miembros de el órgano de representación, en su caso. 
    \item Ratificar las altas acordadas por el órgano de representación y conocer las bajas voluntarias de los asociados.
    \item Resolver, en última instancia, los expedientes relativos a sanción y separación de los asociados, tramitados conforme al procedimiento disciplinario establecido en los presentes Estatutos
    \item Otras que le sean de su competencia en atención a la normativa aplicable. 
\end{enumerate}

\subsection{Adopción de acuerdos}

En las asambleas generales actuarán como presidente/a y secretario/a quienes lo sean de el órgano de representación.

Los acuerdos adoptados por la asamblea general serán recogidos en un acta elaborada y firmada por el secretario del órgano de representación y deberá contar con el visto bueno, mediante su firma, del presidente de la asociación.

Cualquier miembro asistente tiene derecho a solicitar que su intervención o propuestas sean incluidas en el acta.

\subsection{Convocatoria de la asamblea general}

La asamblea general deberá ser convocada al menos en sesión ordinaria una vez al año, dentro del primer trimestre anual, incluyendo, al menos, el siguiente punto del orden del día: examinar y aprobar la liquidación anual de las cuentas del ejercicio anterior, el presupuesto del ejercicio corriente y la memoria de actividades del ejercicio anterior.

Asimismo, se podrá convocar en sesión extraordinaria cuando así lo acuerde el órgano de representación.

Las asambleas generales se celebrarán previa convocatoria del presidente, acompañada del orden del día consignando lugar, fecha y hora. En el caso de que la convocatoria no incluyese el lugar de celebración se entenderá a todos los efectos el domicilio social.

La toma de decisiones se realizará por mayoría de votos de los socios asistentes.

\section{Procedimiento de modificación de estatutos}

Los estatutos de la asociación podrán ser modificados cuando resulte conveniente a los intereses de la misma, por mayoría abasoluta de los votos asistentes en la asamblea general convocada al efecto.

\chapter{Régimen económico}

\section{Patrimonio inicial}

La asociación cuenta con un patrimonio inicial de 66,30 euros.

\section{Ejercicio económico}

El ejercicio económico será anual y su cierre tendrá lugar a 31 de diciembre de cada año.

\section{Recursos económicos}

Constituirán los recursos económicos de la asociación:

\begin{enumerate}
    \item Las cuotas de los miembros, periódicas o extraordinarias.  
    \item Las aportaciones, subvenciones, donaciones a título gratuito, herencias y legados recibidos.  
    \item Bienes muebles e inmuebles.
\end{enumerate}

\section{Cuantía de las cuotas}

Existirán los siguientes tipos de cuotas para los socios:

\begin{itemize}
    \item 20\euro\ al año: precio general.
    \item 10\euro\ al año: precio para estudiantes y/o desempleados.
    \item 45\euro\ al año: precio para ser socios de Python Canarias y Python España.
\end{itemize}

La condición de estudiante y/o desempleado deberá ser justificada documentalmente.

\chapter{Disolución de la asociación}

\section{Causas}

La asociación puede disolverse:

\begin{enumerate}
    \item Por sentencia judicial firme.  
    \item Por acuerdo de la asamblea general extraordinaria.  
    \item Por las causas determinadas en el artículo 39 del Código Civil.
\end{enumerate}

\section{Comisión liquidadora}

Acordada la disolución de la asociación, la asamblea general extraordinaria designará a una comisión liquidadora que tendrá las siguientes funciones:

\begin{enumerate}
    \item Velar por la integridad del patrimonio de la asociación.  
    \item Concluir las operaciones pendientes y efectuar las nuevas, que sean precisas para la liquidación de la asociación.  
    \item Cobrar los créditos de la entidad.  
    \item Liquidar el patrimonio y pagar a los acreedores.  
    \item Aplicar los bienes sobrantes a los fines previstos en los presentes estatutos.  
    \item Solicitar la inscripción de la disolución en el registro de asociaciones.
\end{enumerate}

\chapter*{Disposición final}

Los presentes estatutos han sido aprobados por las siguientes personas fundadoras de la asociación, y de cuyo contenido dan testimonio:

\begin{table}[H]
    \begin{tabular}{l|l}
        {\bf Cargo} & {\bf Nombre y apellidos}\\
        \hline
        Presidente & \textit{Juan Ignacio Rodríguez de León}\\
        Vicepresidente & \textit{Israel Santana Alemán}\\
        Secretario & \textit{Sergio Delgado Quintero}\\
        Tesorero & \textit{Víctor Suárez García}\\
        \hline
        Vocal & \textit{Héctor Álvarez Alonso}\\
        Vocal & \textit{José Lucas Grillo Lorenzo}\\
        Vocal & \textit{Jaime Iván Juanes Prieto}\\
        Vocal & \textit{Víctor Raúl Ruiz Ruiz}\\
        \hline
    \end{tabular}
\end{table}

\begin{flushright}
    En La Orotava, a 1 de octubre de 2018
\end{flushright}

\vspace{15mm}

\begin{multicols}{2}
    \noindent
    El secretario,

    \vspace{25mm}

    \noindent
    Fdo.: \textit{Sergio Delgado Quintero}

    \vfill\null
    \columnbreak
    \begin{flushright}
        El presidente,
    
        \vspace{25mm}
    
        \noindent
        Fdo.: \textit{Juan Ignacio Rodríguez de León}
    \end{flushright}
\end{multicols}


\end{document}
